
\chapter{Related Literature and Studies}
\begin{refsection}

This chapter is divided into three parts: the review of related literature and studies, the synthesis of the state-of-the-art, and the gap bridged by the study. It serves to establish a foundation for the research, deepen the researchers' understanding of the field, and situate the study within the context of existing literature.  


\section{Graph-Based Relational Logic and Perceptual Extraction}

The evolution of automated design frameworks has progressed from basic image matching toward high-precision perceptual analysis. While Shaharom and Tahar [\citeyear{Shaharom2023SIFT}] established the foundational reliability of SIFT and SURF algorithms for maintaining structural integrity and asset preservation during image matching, Garcia and Chen [\citeyear{Garcia2023}] enhanced the precision of these systems by shifting toward the CIE L*a*b* color space for superior perceptual extraction. This transition facilitated a more human-centric approach to identifying brand signatures by separating luminance from chromaticity, thereby effectively reducing visual noise.

Building upon these extraction techniques, Lertrusdachakul, Meesad, and Kittirachakul [\citeyear{Lertrusdachakul2025Hybrid}] addressed the mathematical limitations of earlier clustering models by proposing a Hybrid Generalized K-Means approach to stabilize color centroid determination. Finally, Huang, Chen, and Liu [\citeyear{Huang2025Recoloring}] demonstrated the current versatility of these frameworks by applying intelligent palette mapping to automate complex, high-fidelity design transitions. This progression ensures that modern systems can perform stylistic changes while strictly adhering to the structural "DNA" of the original brand identity.


\section{Graph Neural Networks (GNN) in Spatial Prediction}

The transition toward automated design systems began with the foundational challenge of representing implicit spatial logic through graph-structured data. While Lee, Jiang, Essa, Le, Gong, Yang, and Yang [\citeyear{Lee2020Neural}] introduced the Neural Design Network (NDN) to demonstrate how graph convolution kernels could mimic human spatial reasoning, Nauata, Chang, and Cheng [\citeyear{Nauata2020HouseGAN}] established that relational generative adversarial networks (House-GAN) could encode complex spatial adjacency as bubble diagrams. These early works effectively shifted the field away from rigid, rule-based systems toward models capable of learning the fluid relationships between design elements.

Building on these relational foundations, Nauata, Hosseini, and Chang [\citeyear{Nauata2021HouseGANPlus}] evolved this technology into House-GAN++, a refinement network that moved beyond simple adjacency to demonstrate how GNNs could iteratively transform relational graphs into structurally consistent architectural layouts. By 2024, the application of GNNs expanded into specialized modeling; Jiang, Zhou, Garg, and Oulasvirta [\citeyear{Jiang2024Graph4GUI}] introduced Graph4GUI, utilizing bipartite graphs to achieve high-precision UI element alignment. This was further validated by Purohit, Vyas, and Gupta [\citeyear{Purohit2024Residential}], who confirmed that GNN-based optimization achieves superior performance in constrained residential setups compared to traditional metaheuristic baselines.

Most recently, Kamboj, Kumar, and Deep [\citeyear{Kamboj2025GNN}] culminated this research by demonstrating an end-to-end pipeline in Graphs to Blueprints, capable of converting abstract graph nodes into consistent floor plans at a massive scale. Graph-encoded relational constraints ensure architecturally coherent outputs at scale. This progression—from the initial mimicking of human reasoning to the achievement of high-accuracy UI prediction and large-scale automated modeling—represents the complete evolution of the spatial logic required for modern design orchestration, moving from abstract concepts to a high-precision, mathematically-driven reality.

The development of document image understanding reached a significant turning point with the introduction of LayoutLM by Xu, Li, Cui, Huang, Wei, and Zhou [\citeyear{Xu2020LayoutLM}]. Their research established the methodology of treating text blocks, layout positions, and visual embeddings as semantic nodes, defending the principle that raw page images can be decomposed into meaningful units like paragraphs, figures, and tables. This foundational approach allowed graph-based models to reason over both textual and visual structures simultaneously. By 2022, the field advanced as Han, Wang, Guo, Tang, Wu, and Tian [\citeyear{Han2022VisionGNN}] generalized node extraction beyond documents in their work Vision GNN. By splitting raw images into non-overlapping patches that serve as graph nodes, they demonstrated that patch-based extraction is semantically richer than traditional CNN grids. This methodology enables Graph Neural Networks (GNNs) to capture irregular object boundaries and complex spatial relationships with higher accuracy. Building upon these foundations, Wang, Xu, Cui, and Wei [\citeyear{Wang2023Graphical}] extended node extraction by explicitly modeling text, image, and structural components as heterogeneous nodes within a graph. This progression illustrates a clear trajectory: from the initial establishment of semantic node extraction in 2020, through the extension to general vision tasks in 2022, to the recent culmination in hybrid systems. These modern advancements enable GNNs to reason over complex document and image structures with significantly greater semantic fidelity.

Before a Graph Neural Network can reason over spatial layouts, the raw visual data must first be transformed into discrete, mathematical nodes. Recent literature demonstrates that combining high-speed object detection architectures with Optical Character Recognition (OCR) establishes a robust pipeline for Semantic Layout Parsing. A recent study by dos Santos Júnior, Paixão, and Alvarez [\citeyear{dosSantosJunior2025}], evaluating the comparative performance of YOLOv8 for the layout analysis of document images, concluded that the YOLOv8 architecture provides state-of-the-art precision and speed for extracting structural bounding boxes across complex, heterogeneous layouts. Building upon this validated extraction method, this study passed the generated bounding boxes to the Tesseract OCR engine to retrieve the corresponding textual representations. This hybridized extraction methodology unflattened raster images into the spatial embeddings and labelled text nodes required for downstream GNN processing.

\section{Constraint Satisfaction Problems (CSP) in Design}
The journey toward mathematically maintaining structural integrity during automated layout transformations began with addressing the limitations of traditional linear solvers. Jiang, Stuerzlinger, Zwicker, and Lutteroth [\citeyear{Jiang2020ORCSolver}] developed the ORCSolver, an efficient system designed specifically for adaptive GUI layouts containing non-convex OR-constraints. This breakthrough allowed design frameworks to handle complex scenarios—such as selecting between disparate navigation styles based on viewport space—moving beyond static grids to establish the groundwork for responsive environments that preserve semantic relationships across platforms.

As the demand for more fluid user experiences grew, research shifted toward achieving sub-millisecond response times. Li, Li, Fu, and Cai [\citeyear{Li2024SMT}] introduced SMT-Layout, a MaxSMT-based approach engineered to support real-time interaction in complex interfaces. By formulating layout rules as Satisfiability Modulo Theories (SMT) problems, the system gained the ability to dynamically resolve both soft and hard constraints. This allowed design elements to resize instantly during user interaction, successfully bridging the gap between theoretical constraint solving and practical, high-speed GUI behavior.

The most recent advancements have focused on the dual pillars of scalability and automation. To address the computational demands of modern interfaces, Zhang, Maqrot, Mouysset, and Bortolaso [\citeyear{Zhang2025Scalable}] introduced a CSP-based scalable framework that utilizes decomposition techniques to manage hundreds of interdependent variables without lag. Simultaneously, the manual effort required to define these relationships has been reduced through the work of Penco, Pintar, Vranić, and Šoštarić [\citeyear{Penko2025LLM}], who developed a Large Language Model-driven framework capable of translating natural language requirements into formal constraint models. Together, these developments represent a complete progression from solving specific adaptive logic to enabling the scalable, automated generation of complex, real-time design models.


\section{Stateful Refinement and Iterative Systems}

The development of sophisticated orchestration frameworks begins with the foundational selection of data formats capable of representing complex hierarchical structures. Berdan [\citeyear{Berdan2025Formats}] established the JSON format as a critical standard for lightweight data exchange, identifying its nested structure as uniquely suited for representing the diverse properties of UI nodes. This structural backbone enables seamless communication between frontend interfaces and backend Graph Neural Network (engines), setting the stage for more complex state management. Once this communication channel is established, the focus shifts to maintaining consistency during automated adjustments. Prymushko, Puchko, Yaroshynskyi, Sinko, Kravtsov, and Artemchuk [\citeyear{Prymushko2025CRDT}] addressed this by implementing Conflict-Free Replicated Data Types (CRDTs), which provide a mathematical guarantee for efficient state synchronization in distributed environments. By utilizing CRDTs, a design system can resolve conflicting layout changes without requiring a central coordinator, ensuring the design state remains consistent across all active environments.
 
As synchronization logic became more robust, research moved toward optimizing the speed of these transitions. Gomes, Rodrigues, and Baquero [\citeyear{Gomes2025ConflictSync}] introduced the ConflictSync approach, which utilizes a mechanism for bandwidth-efficient synchronization of divergent states. By focusing on transmitting delta updates rather than full state snapshots, this methodology ensures that complex design graphs can be updated in real-time even across low-bandwidth connections. This efficiency gain proves critical in collaborative environments where multiple agents concurrently modify shared design states.

This technical efficiency is a prerequisite for the final stage of orchestration: the iterative improvement of the system itself. To evaluate this capability, Cheng, Tang, Ma, Hackett, He, Su, Beschastnikh, Huang, Ma, and Xu [\citeyear{Cheng2025SysMoBench}] introduced SysMoBench, a comprehensive benchmark for evaluating AI performance when formally modeling complex real-world systems. Their research emphasizes that refinement is a series of stateful transitions that gradually eliminate design "smells" or constraint violations. By integrating these iterative heuristics with high-speed synchronization, modern frameworks can ensure that every refinement cycle brings the layout closer to the target brand identity while maintaining the structural integrity of the original design intent.

The evolution of bi-directional state machines in user interface research can be traced from early cyber-physical synchronization studies to subsequent applications in human-facing systems. Liang, McGee, Menassa, and Kamat [\citeyear{Liang2020Bidirectional}] established the principle that digital twins and physical robots must maintain continuous reciprocal updates, effectively modeling the interaction as a bi-directional finite state machine. This work highlighted the fundamental necessity of two-way state propagation to ensure total fidelity between simulated models and physical environments.

Building on this foundation, Goldenberg and Tractinsky [\citeyear{Goldenberg2021Bidirectional}] extended the theoretical defense of bi-directionality into the domain of graphical user interfaces. Their study demonstrated that multilingual interfaces—particularly those requiring transitions between right-to-left (RTL) and left-to-right (LTR) orientations—inherently demand two-way directional state transitions. They argued that such interfaces should be formally understood as bi-directional state machines where layout and text orientation dynamically toggle in response to user context. Taken together, these studies illustrate a clear progression: while the foundational work in 2020 defended bi-directional synchronization in cyber-physical systems, subsequent research in 2021 applied the same theoretical framework to multilingual UI design. Subsequent work further extended this framework beyond UI design, exploring how bi-directional state machines could support adaptive accessibility systems. This trajectory grounds bi-directional state machines as a unifying principle that ensures consistency across both machine-to-machine and human-centered interfaces.

\section{Orchestration and Real-Time Processing}

The successful orchestration of automated layout generation within a live, interactive canvas necessitates strict computational performance boundaries. Du, Turner, Dzitsiuk, and Kowdle [\citeyear{Du2020DepthLab}] established that complex geometric operations must be processed within strict millisecond thresholds to prevent user interface occlusion or visual freezing, a principle critical for the viability of algorithmic design frameworks in real-world applications. Furthermore, Hu, Wang, and Wu [\citeyear{Hu2025FacilityLayout}] identified that as generative artificial intelligence paradigms scale to support highly complex spatial layout designs, the computational trade-off between generating optimal configurations and maintaining real-time responsiveness becomes a primary architectural constraint. Consequently, the system must be engineered to resolve mathematical graph dependencies at interactive rates to ensure the orchestration engine does not bottleneck the graphical interface.

To bridge the gap between static algorithmic prediction and dynamic manual refinement, the architecture implements a Human-in-the-Loop (HITL) orchestration paradigm. While traditional end-to-end models often execute procedural planning in a computational vacuum, He, Huang, and Wang [\citeyear{He2022iPLAN}] mathematically proved that computational architectures can co-evolve spatial structures alongside human experts through their iPLAN model. By accommodating free-form interactions such as localized element nudging or resizing, the system treats manual adjustments as definitive, hardened constraints on the canvas. Instead of passively generating a new coordinate array that erases human input, the algorithm proposes procedural suggestions and calculates spatial harmonizations around the user's localized edits. This ensures the orchestration engine remains a collaborative instrument rather than an autonomous override mechanism.

By synthesizing these high-speed constraints and HITL principles, the orchestration layer operated as a highly synchronized, event-driven ecosystem. To prevent computational bottlenecks and ensure the graphical interface did not suffer from visual occlusion or frame-rate degradation during manual manipulation, the framework decoupled continuous drag states from the mathematical solvers. Every manual human refinement acted as an asynchronous, post-interaction event trigger. Specifically, the system hardened the spatial parameters and invoked the underlying constraint solver to calculate optimal coordinate placements only upon the user's release of the element (the on-drop or on-resize-complete event hook). This event-driven approach ensured that the algorithm respected the interactive human inputs without overloading the CPU during fluid drag states, maintaining both global layout harmony and strict latency thresholds \cite{Hu2025FacilityLayout, Du2020DepthLab, He2022iPLAN}.

\section{Algorithmic Evaluation and Aesthetic Scoring}

The evaluation of automated layout architectures necessitates a robust synthesis of high-speed search algorithms and mathematically grounded aesthetic evaluators. Chen, Wang, and Tan [\citeyear{Chen2021Differential}] identified that navigating the nearly infinite permutation space of graphical element placements poses a significant computational bottleneck, as standard linear search models and evolutionary algorithms frequently suffer from latency and stagnation. To resolve these inefficiencies, Vinod and Murthy [\citeyear{Vinod2022Nature}] utilized nature-inspired metaheuristic algorithms to drastically reduce the time-space complexity required for solving engineering layout problems. Specifically, Got, Moussaoui, and Zouache [\citeyear{Got2021Hybrid}] deployed hybrid filter-wrapper mathematical models, such as the Whale Optimization Algorithm, to treat feature selection as a multi-objective task, yielding superior convergence rates. Furthermore, Qais, Mirjalili, and Hasanien [\citeyear{Qais2022Circle}] introduced geometry-based frameworks like the Circle Search Algorithm, which uses circular motion to optimize search trajectories.

Following the generation of spatial coordinates, the framework applies quantitative aesthetic evaluation to score the visual quality of the layouts. Zheng, Yang, and Zeng [\citeyear{Zheng2025Interior}] employed combined CNN-GNN architectures to extract local and global structural features, resulting in an objective compositional scoring system that aligns with strict mathematical rules. Similarly, Guo, Fu, and Zhu [\citeyear{Guo2025Aesthetic}] demonstrated the use of Graph Neural Networks to quantify human-centric design rules—such as whitespace and visual density—into objective metrics. These intelligent design systems utilize AI-driven aesthetic scores to guide evolutionary strategies, successfully automating layouts that balance functionality with visual harmony \cite{zhang2025intelligent}. 


These computational evaluations build upon foundational energy-based layout formulations established by O'Donovan, Agarwala, and Hertzmann [\citeyear{ODonovan2014Layouts}], who demonstrated that graphic layout quality could be mathematically modeled through competing objective functions penalizing whitespace imbalance, boundary collisions, and visual hierarchy degradation. By formulating layout synthesis as an optimization landscape balancing spatial emphasis and negative space, their work established the mathematical basis for evaluating compositional harmony. Expanding this formulation into deep generative models, Zhou, Xu, Ma, Ge, Jiang, and Xu [\citeyear{zhou2022cglgan}] formalized a standardized four-class structural taxonomy for commercial graphic design layouts (comprising logo, text, underlay, and embellishment components) and validated that computational layout generation models must be benchmarked against both objective spatial metrics and structured multi-evaluator designer panels to accurately capture perceptual validity.


Beyond structural alignment and whitespace ratios, modern automated design frameworks must also account for visual interference between generated elements and underlying background content. To prevent the degradation of text legibility caused by visual noise, the layout optimization pipeline of this study integrated dynamic saliency evaluations. Advancing GAN-based domain adaptation for image-aware layout generation, Xu, Zhou, Ge, and Xu [\citeyear{Xuetal2026}] quantified layout occlusion through a spatial clash metric ($R_{shm}$) that computes the perceptual difference in the background's saliency map before and after the application of text masks. Following the same principle, this study measured the background saliency beneath each text element and penalized text placed over visually salient regions (the Legibility Penalty, Formula 3.15). Because the saliency map was computed on each candidate's own composed background, the generated layout preserved the primary focal regions of the background while maintaining structural legibility, without relying on static, offline feature extraction.

To manage the output of real-time generation cycles, the system must prevent the presentation of stochastic or biased candidates. Li, Chen, and Wang [\citeyear{Li2023Fairness}] implemented fairness-aware recommendation ranking algorithms to mathematically balance and sort generated outputs based on rigorous merit. Ultimately, Shi, Shang, and Qi [\citeyear{Shi2023Survey}] emphasized the integration of deep generative models to bridge the gap between high-fidelity coordinate predictions and human-centric aesthetic preferences, a goal that Patnaik, Jain, Krishnamurthy, and Sarkar [\citeyear{Patnaik2025AesthetiQ}] formalized as Aesthetic-Aware Preference Alignment (AAPA). Such alignment provides the necessary heuristic filtering to validate optimal design structures before they are presented to the user.

\section{Synthesis of the State-of-the-Art}

Current state-of-the-art in automated layout generation and graphic design orchestration demonstrates a definitive paradigm shift from manual, coordinate-based drafting toward intelligent, algorithmically-driven frameworks \cite{Kim2025Fragment, Sanchez2025Generative, Hurne2026Blueprint}. Both the established literature and the proposed ThemeSync framework share a foundational reliance on computational algorithms to synthesize spatial arrangements. However, a critical synthesis of the reviewed literature reveals a persistent dichotomy between offline, generative AI models and real-time, interactive design systems. Deep generative models and graph-based frameworks excel at mathematically predicting spatial logic but predominantly operate as static, offline batch processors \cite{Nauata2020HouseGAN, Nauata2021HouseGANPlus}. ThemeSync departs from this passive generation model by embedding advanced computational algorithms into a live, stateful orchestration loop, transforming the canvas from a one-time predictive output into a continuously reactive environment.

This divergence is further evident in the treatment of visual identity and structural logic. Existing literature heavily emphasizes perceptually-conscious clustering models and feature matching algorithms—such as hybrid K-Means, SIFT, and SURF—to isolate dominant hues and stylistic markers with high precision \cite{Garcia2023, Shaharom2023SIFT, Lertrusdachakul2025Hybrid}. Concurrently, Graph Neural Networks (GNNs) represent the current apex of relational spatial modeling, successfully predicting bounding boxes in offline environments \cite{Jiang2024Graph4GUI, Zhou2020Review}. Yet, current frameworks frequently treat these domains in isolation; perceptual extraction is applied primarily to static image analysis, while structural GNNs independently predict layouts. ThemeSync synthesizes these disparate technologies to construct a unified, hybrid "Visual DNA." Current automated design frameworks frequently suffer from typographic hallucinations, as relying purely on Optical Character Recognition (OCR) or feature matching to autonomously extract exact font families from flattened images is computationally error-prone. ThemeSync bridges this gap by decoupling structural extraction from typographic assignment. While the system integrates perceptual extraction directly with relational graph inference to autonomously map colors and stylistic identities onto a live structural graph, the exact typographic parameters—such as font families, specific font weights, and baseline geometric sizing—are deterministically defined by the user prior to orchestration. This ensures that institutional branding remains mathematically consistent during automated transformations, utilizing techniques akin to smart palette mapping \cite{Huang2025Recoloring} to move beyond flat visual adjustments and achieve deep structural awareness.

The most significant technological limitation in existing automated design tools is the prevalence of "destructive updating," where generative models force a complete layout recalculation that irreversibly overwrites manual user refinements. While recent computer science literature validates the necessity of Constraint Satisfaction Problems (CSP) for scalable spatial issues \cite{Li2024SMT, Penko2025LLM, Zhang2025Scalable} and state machines for decentralized synchronization \cite{Liang2020Bidirectional, Prymushko2025CRDT, Gomes2025ConflictSync}, their combined application to protect localized creative intent is notably absent in the design domain. ThemeSync bridges this gap by combining the incremental Cassowary constraint-solving algorithm \cite{Badros2001Cassowary} with a Bi-directional State Machine utilizing JSON-based object serialization \cite{Berdan2025Formats}. Unlike traditional AI generators that act as "fire-and-forget" black boxes, ThemeSync treats user interactions as strict "Hard Constraints." The serialized state machine ensures that these human-in-the-loop refinements survive global theme updates, utilizing the CSP solver to propagate non-destructive recalculations of the surrounding layout and eliminating layout fragility.

Finally, while the literature demonstrates that geometry-based metaheuristics \cite{Got2021Hybrid, Vinod2022Nature, Qais2022Circle} and CNN-GNN architectures \cite{Guo2025Aesthetic, Zheng2025Interior} can mathematically optimize and score layouts based on strict composition rules, finding an optimal layout programmatically does not guarantee subjective user satisfaction. Contemporary generative models frequently force a single, mathematically prioritized output onto the designer. ThemeSync fundamentally contrasts with this approach by synthesizing high-speed state-space exploration with Aesthetic-Aware Preference Alignment (AAPA) \cite{Patnaik2025AesthetiQ}. Rather than automating the designer out of the workflow, the framework utilizes these heuristic-quality filters and fitness functions to evaluate and rank parallel layout candidates in real-time \cite{Li2023Fairness}. This interactive suggestion mechanism maintains the critical "human-in-the-loop" paradigm, proving that high-speed algorithmic orchestration and stateful human creative control can be seamlessly integrated into a single design ecosystem.

\section{Gap Bridged by the Study}

Existing research in graphic design automation has largely advanced along distinct trajectories, with studies either focusing on automated layout generation or supporting manual designer adjustments, but rarely integrating both within a unified framework. As shown in a comprehensive survey on deep generative models \cite{Shi2023Survey}, deep generative models are currently the most popular way to create intelligent layouts. However, these systems work as static, offline pipelines that act like generative black boxes, producing outputs that are structurally correct but cannot be changed by humans during the process. Any changes made by a user are automatically undone when the process is run again. This rigidity is further compounded by the absence of stateful memory in existing design tools, wherein no current system retains a record of designer intent, causing global operations such as theme or template updates to consistently and irreversibly erase localized refinements. Furthermore, even state-of-the-art diffusion-based layout models reduce generation to a purely predictive task, prioritizing mathematical constraints over the subjective aesthetic preferences of designers and failing to align outputs with the contextual demands of professional design practice \cite{Zhang2023LayoutDiffusion, Patnaik2025AesthetiQ}. Taken together, these issues reveal a fundamental and unresolved gap between the operational efficiency of algorithmic automation and its capacity to support meaningful human creative control.

This study fills these gaps by introducing ThemeSync, a stateful graphic design orchestration framework that combines automated layout intelligence with non-destructive human refinement. In response to the offline processing limits identified in prior work \cite{Shi2023Survey}, ThemeSync embeds multi-relational GNN inference and the Cassowary constraint-solving algorithm within a live orchestration loop, allowing structural changes to occur in real time while keeping relational dependencies intact. To resolve the destructive updating problem, the framework employs JSON-based state serialization and state machine logic to encode user interactions as hardened constraints, ensuring localized designer refinements persist through global synchronization events. Furthermore, contemporary deep generative layout models frequently suffer from semantic ambiguity when merging static structural templates with variable data. ThemeSync bridges this gap by decoupling the aesthetic extraction pipeline from a dynamic payload injection system. By utilizing Dynamic Node Instantiation and Implicit Geometric Hierarchy, the framework ensures that user-defined administrative variables—such as specific names or titles—are placed in a defined reading order and assigned role-specific baseline dimensions, allowing dynamic content to be logically orchestrated alongside the master theme. Building on aesthetic alignment principles established in recent literature \cite{Patnaik2025AesthetiQ}, ThemeSync integrates Aesthetic-Aware Preference Alignment (AAPA) and GNN compositional scoring into an interactive recommendation dashboard, replacing single programmatic outputs with a carefully selected set of aesthetically validated layout candidates. ThemeSync changes the role of AI in graphic design from an independent generator to a human-in-the-loop collaborative tool that respects both structural logic and creative intent.

% \begin{rotatepage}
% \begin{sidewaysfigure}[!t]
%     \centering
%     \includegraphics[width=\textwidth]{figures/sampleFig2.png}
%     \caption[Black-body radiation ddd dddd dddd]{Spectra of black-body radiation at various temperatures, according to Wien's displacement law \cite{wannier1987statistical}.}
%     \label{fig:secondFig}
% \end{sidewaysfigure}
% \end{rotatepage}

% \begin{figure}[ht]
%     \centering
% 	\includegraphics[width=0.85\textwidth]{figures/sampleFig1.jpg} 
% 	\caption[Barred spiral galaxy NGC 1300]{Barred spiral galaxy NGC 1300 photographed by Hubble telescope. While the galaxy in the photo is not our sun, it does emit light, much like our sun. Image credit: NASA.}
% 	\label{fig:firstFig}
% \end{figure}



%=======================================================%
%%%%% Do not delete this part %%%%%%
\clearpage

\begin{center}
\textbf{Notes}
\end{center}
\printbibliography[heading=none]
\end{refsection}